% ============================================================
%  Übungsaufgaben Template (my_dhbw Stil, uebung_*.tex)
%  Header "Übungsblatt", grün eingefärbte <details>-Lösungen.
%  Renderer: pdflatex (zweimal)
% ============================================================
\documentclass[11pt, a4paper]{article}

\usepackage[ngerman]{babel}
\usepackage{fontspec}
\setmainfont[
  Path=/usr/share/fonts/TTF/,
  Extension=.ttf,
  BoldFont=DejaVuSans-Bold,
  ItalicFont=DejaVuSans-Oblique,
  BoldItalicFont=DejaVuSans-BoldOblique
]{DejaVuSans}
\setsansfont[
  Path=/usr/share/fonts/TTF/,
  Extension=.ttf,
  BoldFont=DejaVuSans-Bold,
  ItalicFont=DejaVuSans-Oblique,
  BoldItalicFont=DejaVuSans-BoldOblique
]{DejaVuSans}
\setmonofont[Path=/usr/share/fonts/TTF/,Extension=.ttf]{DejaVuSansMono}
\usepackage[top=2cm, bottom=2cm, left=2.4cm, right=2.4cm, footskip=1cm]{geometry}
\usepackage{amsmath, amssymb}
\usepackage{xcolor}
\usepackage{enumitem}
\usepackage{fancyhdr}
\usepackage{lastpage}
\usepackage{tabularx}
\usepackage{booktabs}
\usepackage{microtype}
\usepackage{parskip}

\usepackage{array}
\usepackage{longtable}
\usepackage{ltablex}
\keepXColumns
\usepackage{colortbl}
\usepackage[most,breakable]{tcolorbox}
\usepackage{listings}
\usepackage[normalem]{ulem}
\usepackage{pdflscape}
\usepackage{titlesec}
\usepackage[colorlinks=true, linkcolor=accent, urlcolor=accent]{hyperref}

\renewcommand{\familydefault}{\sfdefault}

% ── Theme ─────────────────────────────────────────────────────
\definecolor{accent}{RGB}{0, 60, 113}
\definecolor{primaryblue}{RGB}{0, 60, 113}
\definecolor{loesung}{RGB}{0, 100, 0}
\definecolor{codebg}{RGB}{245, 245, 245}
\definecolor{figurebg}{RGB}{232, 241, 255}
\definecolor{tablehintbg}{RGB}{240, 255, 240}
\definecolor{solutionbg}{RGB}{240, 252, 240}
\definecolor{rulegray}{RGB}{180, 180, 180}
\definecolor{tableheadbg}{RGB}{0, 60, 113}

% ── Header/Footer ─────────────────────────────────────────────
\pagestyle{fancy}
\fancyhf{}
\renewcommand{\headrulewidth}{0.4pt}
\fancyhead[L]{\footnotesize\color{accent}\nichttitle}
\fancyhead[R]{\footnotesize\color{accent}\nichtsubtitle}
\fancyfoot[R]{\footnotesize\color{accent}Blatt \thepage\,/\,\pageref{LastPage}}

% ── Typografie ────────────────────────────────────────────────
\titleformat{\section}
    {\color{accent}\large\bfseries}{}{0pt}{}
    [\vspace{-4pt}\color{accent}\rule{\linewidth}{1.5pt}]
\titleformat{\subsection}
    {\normalsize\bfseries\color{accent!85!black}}{}{0pt}{}{}
\titleformat{\subsubsection}
    {\small\bfseries\color{accent!70!black}}{}{0pt}{}{}
\titlespacing*{\section}{0pt}{10pt}{3pt}
\titlespacing*{\subsection}{0pt}{6pt}{2pt}
\titlespacing*{\subsubsection}{0pt}{4pt}{1pt}

\setlength{\parindent}{0pt}
\setlength{\parskip}{6pt}
\renewcommand{\arraystretch}{1.2}
\setlist[itemize]{noitemsep, topsep=3pt, parsep=0pt, leftmargin=1.4em}
\setlist[enumerate]{noitemsep, topsep=3pt, parsep=0pt, leftmargin=1.6em}

% ── Boxen: solutionbox = Lösungsbox in grün ──────────────────
\newtcolorbox{figurebox}[1]{
  colback=figurebg, colframe=accent!50,
  title={Abbildung: #1}, fonttitle=\small\bfseries,
  breakable, before skip=6pt, after skip=6pt
}
\newtcolorbox{tablehintbox}[1]{
  colback=tablehintbg, colframe=green!50!black!50,
  title={Tabelle: #1}, fonttitle=\small\bfseries,
  breakable, before skip=6pt, after skip=6pt
}
\newtcolorbox{solutionbox}[1]{
  enhanced,
  colback=solutionbg, colframe=loesung,
  title={#1}, fonttitle=\small\bfseries,
  coltitle=white, colbacktitle=loesung,
  breakable, before skip=8pt, after skip=8pt
}

\lstset{
  basicstyle=\ttfamily\small,
  backgroundcolor=\color{codebg},
  breaklines=true,
  frame=single, framerule=0pt,
  xleftmargin=4pt, xrightmargin=4pt,
  aboveskip=6pt, belowskip=6pt,
  keepspaces=true
}

\providecommand{\nichttitle}{Übungsblatt}
\providecommand{\nichtsubtitle}{}

\renewcommand{\nichttitle}{Relationen, Algebra, Diskrete Mathematik}
\renewcommand{\nichtsubtitle}{Übungsblatt 9}


\begin{document}

\begin{center}
{\Large\bfseries\color{accent}\nichttitle}
\ifx\nichtsubtitle\empty\else
  \\[2pt]{\normalsize\color{accent!80!black}\nichtsubtitle}
\fi
\\[2pt]
\color{accent}\rule{0.4\linewidth}{0.8pt}\color{black}
\end{center}
\vspace{4pt}

% ── BODY ──────────────────────────────────────────────────────

\section{Übungsblatt 9 — Graphen}


\noindent\textcolor{rulegray}{\rule{\linewidth}{0.4pt}}


\paragraph{Aufgabe 1 — Begriffe und Knotengrade \textit{(8 Punkte)}}\mbox{}\\[2pt]

a) Definiere die Begriffe \textbf{Knotengrad}, \textbf{Multigraph} und \textbf{bipartiter Graph} jeweils in einem Satz. \textit{(3 P)}


b) Gegeben sei der ungerichtete Graph G = (V, E) mit V = \{1, 2, 3, 4, 5\} und

E = \{\{1,2\}, \{1,3\}, \{2,3\}, \{2,4\}, \{3,4\}, \{4,5\}\}.

Bestimme den Knotengrad jedes Knotens und die Summe aller Knotengrade. Begründe, warum diese Summe stets gerade sein muss. \textit{(5 P)}


\textbf{Lösung:}


a)

\begin{itemize}
  \item \textbf{Knotengrad}: Anzahl der an einen Knoten angrenzenden Kanten.
  \item \textbf{Multigraph}: Graph, in dem zwischen zwei Knoten mehrere (parallele) Kanten erlaubt sind.
  \item \textbf{Bipartiter Graph}: Graph, dessen Knotenmenge in zwei disjunkte Mengen A und B zerlegt werden kann, sodass jede Kante einen Knoten aus A mit einem Knoten aus B verbindet.
\end{itemize}

b) Knotengrade:

\begin{itemize}
  \item deg(1) = 2 (Kanten zu 2, 3)
  \item deg(2) = 3 (Kanten zu 1, 3, 4)
  \item deg(3) = 3 (Kanten zu 1, 2, 4)
  \item deg(4) = 3 (Kanten zu 2, 3, 5)
  \item deg(5) = 1 (Kante zu 4)
\end{itemize}

Summe = 2 + 3 + 3 + 3 + 1 = \textbf{12}.


Begründung: Jede Kante \{u, v\} trägt zu deg(u) und zu deg(v) jeweils 1 bei, also insgesamt 2. Die Gradsumme ist daher gleich 2·|E| und damit immer gerade (Handschlaglemma). Hier: 2·6 = 12. ✓



\noindent\textcolor{rulegray}{\rule{\linewidth}{0.4pt}}


\paragraph{Aufgabe 2 — Darstellungsformen umwandeln \textit{(15 Punkte)}}\mbox{}\\[2pt]

Ein gerichteter, gewichteter Graph beschreibt Flugverbindungen mit Flugzeit (in Stunden) zwischen vier Städten:

A = Stuttgart, B = Berlin, C = Wien, D = Paris.


Verbindungen: A→B (1.0), A→D (1.5), B→C (1.0), C→A (1.0), C→D (2.0), D→B (1.5).


a) Stelle den Graphen als \textbf{Adjazenzmatrix} auf (Zeilenreihenfolge A, B, C, D; nicht vorhandene Kanten als 0). \textit{(5 P)}


b) Gib die \textbf{Adjazenzliste} mit jeweils Vorgänger- und Nachfolgerlisten an. \textit{(6 P)}


c) Beurteile, welche Darstellung für diesen konkreten Graphen speichereffizienter ist und ab welcher Kantendichte die andere Darstellung vorteilhafter würde. \textit{(4 P)}


\textbf{Lösung:}


a) Adjazenzmatrix A (Zeile = Start, Spalte = Ziel):



{\small
\begin{tabularx}{\linewidth}{XXXXX}
\hline
\rowcolor{tableheadbg}
{\color{white}\textbf{}} & {\color{white}\textbf{A}} & {\color{white}\textbf{B}} & {\color{white}\textbf{C}} & {\color{white}\textbf{D}} \\[2pt]
\hline
\endhead
\hline
\endfoot
A & 0 & 1.0 & 0 & 1.5 \\
B & 0 & 0 & 1.0 & 0 \\
C & 1.0 & 0 & 0 & 2.0 \\
D & 0 & 1.5 & 0 & 0 \\
\end{tabularx}
}


Die Matrix ist nicht symmetrisch — der Graph ist gerichtet.


b) Adjazenzlisten (gerichtet → zwei Listen pro Knoten):



{\small
\begin{tabularx}{\linewidth}{XXX}
\hline
\rowcolor{tableheadbg}
{\color{white}\textbf{Knoten}} & {\color{white}\textbf{Nachfolger (Out)}} & {\color{white}\textbf{Vorgänger (In)}} \\[2pt]
\hline
\endhead
\hline
\endfoot
A & B(1.0), D(1.5) & C(1.0) \\
B & C(1.0) & A(1.0), D(1.5) \\
C & A(1.0), D(2.0) & B(1.0) \\
D & B(1.5) & A(1.5), C(2.0) \\
\end{tabularx}
}


c) Speicherbedarf:

\begin{itemize}
  \item Adjazenzmatrix: n² = 16 Einträge, davon nur 6 belegt → 10 Nullen werden mitgespeichert.
  \item Adjazenzliste: 6 Kanten × (Ziel + Gewicht) ≈ 12 Einträge plus Listenverwaltung.
\end{itemize}

Der Graph ist \textbf{spärlich} (|E| = 6, max. n(n−1) = 12 Kanten, also ca. 50\% — bei realen Flugnetzen meist deutlich weniger). Die Adjazenzliste skaliert mit O(|V| + |E|), die Matrix mit O(|V|²). Für dichte Graphen mit |E| ≈ |V|² (etwa ab 60–70 \% Dichte) wird die Matrix vorteilhaft, weil dann die Listenverwaltung den Vorteil verliert und der Matrix-Zugriff in O(1) auf einzelne Kanten möglich ist.



\noindent\textcolor{rulegray}{\rule{\linewidth}{0.4pt}}


\paragraph{Aufgabe 3 — Graphtypen identifizieren und konstruieren \textit{(12 Punkte)}}\mbox{}\\[2pt]

a) Sei G ein ungerichteter Graph mit |V| = 6 und |E| = 6, in dem jeder Knoten Grad 2 hat. Welcher \textbf{Graphtyp aus der Tabelle} ist G zwingend, und welcher nicht? Begründe. \textit{(4 P)}


b) Konstruiere einen \textbf{vollständigen} ungerichteten Graphen K₅. Wie viele Kanten hat er allgemein für n Knoten? Wende die Formel auf n = 5 und n = 8 an. \textit{(4 P)}


c) Zeichne (oder beschreibe per Kantenliste) einen \textbf{bipartiten} Graphen mit den Mengen X = \{x₁, x₂, x₃\} und Y = \{y₁, y₂\} so, dass jeder Knoten aus X mit genau zwei Knoten aus Y und jeder Knoten aus Y mit genau drei Knoten aus X verbunden ist. Prüfe die Konsistenz der Gradsummen. \textit{(4 P)}


\textbf{Lösung:}


a) Es gilt |E| = |V| = 6 und alle Knotengrade = 2 — das ist exakt die Definition eines \textbf{Kreises} aus der Typen-Tabelle. Es ist allerdings nicht zwingend ein \textbf{einziger} Kreis: G könnte auch in zwei disjunkte Dreiecke (3-Kreise) zerfallen. Ein \textbf{Baum} ist G nicht (Bäume sind zyklenfrei, hier gibt es Zyklen), ebenso wenig ein \textbf{Stern} (dort hätte ein Zentrum Grad n−1 = 5 und alle anderen Grad 1).


b) Im vollständigen Graphen K\_n verbindet jede Kante zwei verschiedene Knoten, jedes Paar genau einmal:

\[
|E| = \binom{n}{2} = \frac{n(n-1)}{2}
\]


\begin{itemize}
  \item n = 5: |E| = 5·4/2 = \textbf{10} Kanten.
  \item n = 8: |E| = 8·7/2 = \textbf{28} Kanten.
\end{itemize}

K₅ als Kantenliste: \{1-2, 1-3, 1-4, 1-5, 2-3, 2-4, 2-5, 3-4, 3-5, 4-5\}.


c) Geforderte Grade: deg(xᵢ) = 2 für i = 1,2,3 und deg(yⱼ) = 3 für j = 1,2.


Konsistenzprüfung über die Gradsummen:

\begin{itemize}
  \item Summe in X: 3·2 = 6.
  \item Summe in Y: 2·3 = 6.
\end{itemize}

Beide Summen müssen gleich sein, weil jede Kante eines bipartiten Graphen genau einen Endpunkt in X und einen in Y hat. ✓


Eine mögliche Kantenmenge:

E = \{\{x₁,y₁\}, \{x₁,y₂\}, \{x₂,y₁\}, \{x₂,y₂\}, \{x₃,y₁\}, \{x₃,y₂\}\}.


Das ist gerade der vollständige bipartite Graph K\_\{3,2\} mit 6 Kanten.



\noindent\textcolor{rulegray}{\rule{\linewidth}{0.4pt}}


\paragraph{Aufgabe 4 — Eulersche und Hamiltonsche Pfade \textit{(18 Punkte)}}\mbox{}\\[2pt]

Gegeben sei der ungerichtete Graph H = (V, E) mit V = \{1, 2, 3, 4, 5\} und folgender Adjazenzmatrix (1 = Kante, 0 = keine Kante):



{\small
\begin{tabularx}{\linewidth}{XXXXXX}
\hline
\rowcolor{tableheadbg}
{\color{white}\textbf{}} & {\color{white}\textbf{1}} & {\color{white}\textbf{2}} & {\color{white}\textbf{3}} & {\color{white}\textbf{4}} & {\color{white}\textbf{5}} \\[2pt]
\hline
\endhead
\hline
\endfoot
1 & 0 & 1 & 1 & 1 & 0 \\
2 & 1 & 0 & 1 & 0 & 1 \\
3 & 1 & 1 & 0 & 1 & 1 \\
4 & 1 & 0 & 1 & 0 & 1 \\
5 & 0 & 1 & 1 & 1 & 0 \\
\end{tabularx}
}


a) Bestimme die Knotengrade und entscheide anhand der Grade, ob ein \textbf{Eulerscher Zyklus} existieren kann. Bestimme zusätzlich, ob ein \textbf{Eulerscher Pfad} (offen, Anfang ≠ Ende) existiert. \textit{(6 P)}


b) Gib einen \textbf{Hamiltonschen Zyklus} in H an oder begründe, warum keiner existiert. \textit{(6 P)}


c) Erläutere den \textbf{Unterschied} zwischen Eulerschem und Hamiltonschem Pfad in eigenen Worten und gib ein Beispiel für einen Graphen an, der einen Eulerschen Zyklus besitzt, aber keinen Hamiltonschen Zyklus. \textit{(6 P)}


\textbf{Lösung:}


a) Knotengrade aus der Zeilensumme:

\begin{itemize}
  \item deg(1) = 3, deg(2) = 3, deg(3) = 4, deg(4) = 3, deg(5) = 3.
\end{itemize}

Gesamtkantenzahl = (3+3+4+3+3)/2 = 8.


Eulerscher Zyklus: existiert genau dann, wenn der Graph zusammenhängend ist und \textbf{alle} Knoten \textbf{geraden Grad} haben. Hier haben vier Knoten Grad 3 (ungerade) → \textbf{kein Eulerscher Zyklus}.


Eulerscher (offener) Pfad: existiert genau dann, wenn der Graph zusammenhängend ist und \textbf{genau zwei} Knoten ungeraden Grad haben (das sind dann Anfang und Ende). Hier sind es \textbf{vier} ungerade Knoten → \textbf{kein Eulerscher Pfad}.


b) Ein Hamiltonscher Zyklus muss jeden Knoten genau einmal besuchen und am Start enden. Versuch:


1 → 2 → 5 → 4 → 3 → 1.


Prüfung der Kanten anhand der Matrix: \{1,2\}✓, \{2,5\}✓, \{5,4\}✓, \{4,3\}✓, \{3,1\}✓. Alle fünf Knoten kommen je einmal vor, der Zyklus schließt bei 1. → \textbf{Hamiltonscher Zyklus existiert}: (1, 2, 5, 4, 3, 1).


c) \textbf{Unterschied}: Ein Eulerscher Pfad benutzt \textbf{jede Kante genau einmal} (Knoten dürfen mehrfach besucht werden). Ein Hamiltonscher Pfad besucht \textbf{jeden Knoten genau einmal} (Kanten dürfen also nicht alle benutzt werden). Eulerschheit ist eine Eigenschaft über die Kantenmenge und effizient (über Knotengrade) entscheidbar; Hamiltonschheit ist eine Eigenschaft über die Knotenmenge und im Allgemeinen NP-schwer zu entscheiden.


\textbf{Beispiel}: Zwei Dreiecke, die sich an einem gemeinsamen Knoten v berühren (Knotenmenge \{a, b, v, c, d\}, Kanten \{a-b, a-v, b-v, c-v, d-v, c-d\}). Alle Knotengrade sind gerade (a, b, c, d je 2; v hat 4) → Eulerscher Zyklus existiert. Ein Hamiltonscher Zyklus müsste den „Schnittknoten" v aber zweimal nutzen, um beide Dreiecke abzulaufen — das verstößt gegen die Hamilton-Bedingung. Also kein Hamiltonscher Zyklus.



\noindent\textcolor{rulegray}{\rule{\linewidth}{0.4pt}}


\paragraph{Aufgabe 5 — Fehleranalyse einer Adjazenzmatrix \textit{(10 Punkte)}}\mbox{}\\[2pt]

Eine Studentin behauptet, folgende Adjazenzmatrix repräsentiere einen \textbf{ungerichteten, einfachen} (kein Multigraph, keine Schleifen) Graphen mit 4 Knoten:



{\small
\begin{tabularx}{\linewidth}{XXXXX}
\hline
\rowcolor{tableheadbg}
{\color{white}\textbf{}} & {\color{white}\textbf{1}} & {\color{white}\textbf{2}} & {\color{white}\textbf{3}} & {\color{white}\textbf{4}} \\[2pt]
\hline
\endhead
\hline
\endfoot
1 & 0 & 1 & 2 & 0 \\
2 & 1 & 1 & 1 & 0 \\
3 & 1 & 1 & 0 & 1 \\
4 & 0 & 0 & 1 & 0 \\
\end{tabularx}
}


a) Identifiziere \textbf{alle Verstöße} gegen die Eigenschaften eines ungerichteten, einfachen Graphen und benenne sie konkret mit Position in der Matrix. \textit{(6 P)}


b) Korrigiere die Matrix minimal (möglichst wenige Änderungen), sodass sie einen ungerichteten, einfachen Graphen darstellt. Gib die korrigierte Matrix und die zugehörige Kantenmenge an. \textit{(4 P)}


\textbf{Lösung:}


a) Drei Verstöße:


\begin{enumerate}
  \item \textbf{Nicht symmetrisch}: a₁₃ = 2, aber a₃₁ = 1. Bei einem ungerichteten Graphen muss A = Aᵀ gelten, weil eine Kante \{i, j\} in beiden Richtungen identisch eingetragen wird.
  \item \textbf{Multigraph-Eintrag}: a₁₃ = 2 deutet bei einem Multigraphen auf zwei parallele Kanten zwischen 1 und 3 hin. In einem einfachen Graphen sind nur Werte 0 und 1 zulässig.
  \item \textbf{Schleife}: a₂₂ = 1 bedeutet eine Kante von Knoten 2 zu sich selbst. Schleifen sind in einfachen Graphen nicht erlaubt.
\end{enumerate}

b) Minimale Korrektur:

\begin{itemize}
  \item a₁₃ = a₃₁ = 1 (Symmetrie herstellen, Mehrfachkante auf einfache Kante reduzieren).
  \item a₂₂ = 0 (Schleife entfernen).
\end{itemize}

Korrigierte Matrix:



{\small
\begin{tabularx}{\linewidth}{XXXXX}
\hline
\rowcolor{tableheadbg}
{\color{white}\textbf{}} & {\color{white}\textbf{1}} & {\color{white}\textbf{2}} & {\color{white}\textbf{3}} & {\color{white}\textbf{4}} \\[2pt]
\hline
\endhead
\hline
\endfoot
1 & 0 & 1 & 1 & 0 \\
2 & 1 & 0 & 1 & 0 \\
3 & 1 & 1 & 0 & 1 \\
4 & 0 & 0 & 1 & 0 \\
\end{tabularx}
}


Sie ist symmetrisch, enthält nur 0/1, und die Diagonale ist Null. Kantenmenge: E = \{\{1,2\}, \{1,3\}, \{2,3\}, \{3,4\}\}. Knotengrade: deg(1) = 2, deg(2) = 2, deg(3) = 3, deg(4) = 1; Gradsumme 8 = 2·|E| ✓.



\noindent\textcolor{rulegray}{\rule{\linewidth}{0.4pt}}


\paragraph{Aufgabe 6 — Transfer: Modellierung eines Lerngruppen-Netzwerks \textit{(17 Punkte)}}\mbox{}\\[2pt]

In einem DHBW-Kurs gibt es 6 Studierende: A, B, C, D, E, F. Eine Plattform protokolliert, wer mit wem \textbf{gemeinsam an Aufgaben arbeitet} (ungerichtete Beziehung) und zusätzlich, \textbf{wer wem regelmäßig Material zuschickt} (gerichtete Beziehung). Aus den Logs ergibt sich:


\begin{itemize}
  \item Gemeinsames Arbeiten: \{A,B\}, \{A,C\}, \{B,C\}, \{C,D\}, \{D,E\}, \{E,F\}, \{D,F\}.
  \item Material-Zusenden (Sender → Empfänger): A→C, C→A, C→B, D→C, E→D, F→D, F→E.
\end{itemize}

a) Modelliere das gemeinsame Arbeiten als ungerichteten Graphen G₁ und ordne ihn einem Typ aus der Tabelle zu, sofern möglich. Gib die Knotengrade an. \textit{(5 P)}


b) Modelliere das Material-Zusenden als gerichteten Graphen G₂. Bestimme zu jedem Knoten \textbf{Eingangsgrad} (Anzahl eingehender Kanten) und \textbf{Ausgangsgrad} (Anzahl ausgehender Kanten). \textit{(5 P)}


c) Die Plattform soll erkennen, ob es im Graphen G₂ einen \textbf{Zyklus} gibt (also eine Gruppe, die im Kreis Material austauscht). Gib alle Zyklen in G₂ an und erläutere den Unterschied zwischen einem \textbf{Pfad} und einem \textbf{Zyklus} anhand eines konkreten Beispiels aus deinem Graphen. \textit{(7 P)}


\textbf{Lösung:}


a) G₁ = (V, E₁) mit V = \{A, B, C, D, E, F\} und

E₁ = \{\{A,B\}, \{A,C\}, \{B,C\}, \{C,D\}, \{D,E\}, \{E,F\}, \{D,F\}\}, also |E₁| = 7.


Knotengrade:

\begin{itemize}
  \item deg(A) = 2 (B, C)
  \item deg(B) = 2 (A, C)
  \item deg(C) = 3 (A, B, D)
  \item deg(D) = 3 (C, E, F)
  \item deg(E) = 2 (D, F)
  \item deg(F) = 2 (D, E)
\end{itemize}

Gradsumme = 14 = 2·7 ✓.


Typzuordnung: G₁ ist \textbf{kein} Baum (es gibt Zyklen, z. B. A-B-C-A). Er ist \textbf{kein Kreis} (|E| = 7 ≠ |V| = 6, und nicht alle Grade sind 2). Er ist \textbf{kein Stern} (kein Zentrum von Grad 5). Er ist \textbf{kein vollständiger Graph} (z. B. fehlt \{A,D\}). Er ist auch \textbf{nicht bipartit} (das Dreieck A-B-C ist ein ungerader Zyklus). Er entspricht damit \textbf{keinem} der genannten Spezialtypen — es ist einfach ein zusammenhängender ungerichteter Graph, der aus zwei Dreiecken (A-B-C und D-E-F) besteht, die durch die Brücke C-D verbunden sind.


b) G₂ = (V, E₂) gerichtet mit

E₂ = \{(A,C), (C,A), (C,B), (D,C), (E,D), (F,D), (F,E)\}, |E₂| = 7.



{\small
\begin{tabularx}{\linewidth}{XXX}
\hline
\rowcolor{tableheadbg}
{\color{white}\textbf{Knoten}} & {\color{white}\textbf{Eingangsgrad (in)}} & {\color{white}\textbf{Ausgangsgrad (out)}} \\[2pt]
\hline
\endhead
\hline
\endfoot
A & 1 (von C) & 1 (nach C) \\
B & 1 (von C) & 0 \\
C & 2 (A, D) & 2 (A, B) \\
D & 2 (E, F) & 1 (nach C) \\
E & 1 (von F) & 1 (nach D) \\
F & 0 & 2 (D, E) \\
\end{tabularx}
}


Konsistenzprüfung: Σ in = Σ out = 7 = |E₂| ✓ (in gerichteten Graphen gilt Σ in-deg = Σ out-deg = |E|).


c) \textbf{Zyklen in G₂}: Ein Zyklus ist ein geschlossener Pfad mit unterschiedlichen Kanten, der zum Startknoten zurückführt.


\begin{itemize}
  \item (A, C, A): A→C→A ist ein 2-Zyklus über die Kanten (A,C) und (C,A). Die beiden gerichteten Kanten sind unterschiedlich, also ist es ein gültiger Zyklus.
  \item Weitere Zyklen existieren nicht: F hat in-deg 0 und kann daher in keinem Zyklus liegen; E geht nur an D, D nur an C, C verzweigt zu A oder B; B ist ein Senken-Knoten (out-deg 0). Damit ist die einzige geschlossene Struktur das Paar A↔C.
\end{itemize}

\textbf{Pfad vs. Zyklus} (konkretes Beispiel):

\begin{itemize}
  \item Der Pfad F → E → D → C → B besucht jeden Knoten höchstens einmal, beginnt bei F und endet bei B (Anfang ≠ Ende, offen).
  \item Der Zyklus A → C → A schließt sich (Anfang = Ende), und die durchlaufenen Kanten (A,C) und (C,A) sind verschieden.
\end{itemize}

Kurz: Ein Pfad ist eine Knoten-/Kantensequenz ohne Knotenwiederholung; ein Zyklus ist ein geschlossener Pfad, dessen Endknoten gleich dem Startknoten ist und dessen Kanten alle verschieden sind. Bäume aus der Typen-Tabelle sind genau die zusammenhängenden Graphen ohne jeden Zyklus — G₂ ist also kein gerichteter Baum.





% ──────────────────────────────────────────────────────────────

\end{document}
